\section{Mesin Abstrak dan Daya Komputasi}
\label{sec:ch02-mesin}

\subsection{Otomata sebagai Pengenal Bahasa}

Otomata adalah model matematika dari mesin yang membaca string simbol
satu per satu dan berpindah antarstate. Secara informal, otomata
memiliki himpunan state, state awal, himpunan state penerima, dan
fungsi transisi. String diterima jika setelah seluruh simbol dibaca,
mesin berada pada state penerima.

Otomata hingga (\eng{finite automaton}) memiliki memori terbatas:
hanya state saat ini yang disimpan. Karena itu, otomata hingga hanya
mengenali bahasa regular. Otomata pushdown menambahkan stack, sehingga
dapat mengenali bahasa bebas konteks. Mesin Turing menambahkan pita
yang dapat dibaca dan ditulis tanpa batas, sehingga dapat mengenali
bahasa yang lebih luas.

\begin{definisi}[DFA]
Deterministic Finite Automaton (DFA) adalah lima-tupel
\[
M = (Q, \Sigma, \delta, q_0, F),
\]
dengan \(Q\) himpunan state berhingga, \(\Sigma\) alfabet, \(q_0\)
state awal, \(F \subseteq Q\) himpunan state penerima, dan
\(\delta : Q \times \Sigma \to Q\) fungsi transisi total.
\end{definisi}

DFA bersifat deterministik: untuk setiap state dan simbol, terdapat
tepat satu state tujuan. NFA (\eng{nondeterministic finite automaton})
mengizinkan nol, satu, atau beberapa transisi untuk pasangan yang sama,
termasuk transisi \(\varepsilon\). Meskipun tampak lebih kuat, NFA dan
DFA mengenali kelas bahasa yang sama; buktinya diberikan melalui
konstruksi subset pada Bab VI.

\subsection{Hierarki Chomsky}

\begin{table}[H]
\centering
\caption{Hierarki Chomsky dan mesin pengenal}
\label{tab:ch02-chomsky}
\begin{tabularx}{\textwidth}{@{}l l X l@{}}
\toprule
\textbf{Tipe} & \textbf{Kelas bahasa} & \textbf{Bentuk aturan} &
\textbf{Mesin} \\
\midrule
3 & Regular & \(A \to aB\), \(A \to a\), \(A \to \varepsilon\) &
Otomata hingga \\
2 & Bebas konteks & \(A \to \alpha\) & Otomata pushdown \\
1 & Sensitif konteks & \(\alpha A \beta \to \alpha \gamma \beta\) &
Otomata terbatas linear \\
0 & Rekursif terbilang & \(\alpha \to \beta\) & Mesin Turing \\
\bottomrule
\end{tabularx}
\end{table}

Setiap kelas bahasa merupakan himpunan bagian sejati dari kelas di
atasnya. Bahasa \(\{a^n b^n \mid n \ge 0\}\) adalah contoh bahasa
bebas konteks yang bukan regular. Bahasa
\(\{a^n b^n c^n \mid n \ge 0\}\) adalah contoh bahasa sensitif konteks
yang bukan bebas konteks.

\begin{figure}[H]
\centering
\begin{tikzpicture}[node distance=4mm]
  \node[phase, text width=4.2cm] (reg) {Tipe 3: Regular\\Otomata hingga};
  \node[phase, text width=4.2cm, below=of reg] (cfg) {Tipe 2: Bebas konteks\\Otomata pushdown};
  \node[phase, text width=4.2cm, below=of cfg] (cs) {Tipe 1: Sensitif konteks\\Otomata terbatas linear};
  \node[phase, text width=4.2cm, below=of cs] (re) {Tipe 0: Rekursif terbilang\\Mesin Turing};
  \draw[arrow] (reg) -- (cfg);
  \draw[arrow] (cfg) -- (cs);
  \draw[arrow] (cs) -- (re);
\end{tikzpicture}
\caption{Tangga hierarki Chomsky: semakin ke bawah, semakin besar daya ekspresif}
\label{fig:ch02-chomsky}
\end{figure}

\subsection{Implikasi bagi Kompilator}

Front-end kompilator memanfaatkan dua tingkat hierarki. Analisis
leksikal berada pada Tipe 3 karena token dapat dikenali dengan memori
terbatas. Analisis sintaksis berada pada Tipe 2 karena struktur
program membutuhkan stack, misalnya untuk menangani kurung bersarang
atau pemanggilan fungsi. Analisis semantik melampaui Tipe 2 karena
memerlukan informasi konteks seperti tipe dan scope.

Pemahaman hierarki ini mencegah kesalahan Desain: ekspresi reguler
tidak cukup untuk memeriksa kurung seimbang, dan tata bahasa bebas
konteks tidak cukup untuk memeriksa kesesuaian tipe. Setiap tingkat
memerlukan alat dan teknik yang berbeda.
